\documentclass[10pt,a4paper]{amsart}
\usepackage[margin=19mm]{geometry}
\usepackage{amsmath,amssymb,mathtools}
\usepackage[scr=rsfs,cal=euler]{mathalfa}
\usepackage{xfrac}
\usepackage[unicode,hidelinks,bookmarks=false]{hyperref}
\newcommand{\ie}{i.e.\ }
\newcommand{\cf}{cf.\ }
\newcommand{\resp}{respectively\ }
\newcommand{\Subsection}{\subsection}
\providecommand{\nobreakdash}{\nobreak\hbox{-}}
\setlength{\emergencystretch}{2em}
\multlinegap0pt
\usepackage{bm}
\usepackage{bbm}
\usepackage{dsfont}
\usepackage{xspace}
\usepackage{cancel}
\usepackage{mathtools}
\usepackage{amsthm}
\makeatletter
\@ifundefined{theorem}{\newtheorem{theorem}{Theorem}}{}
\@ifundefined{lemma}{\newtheorem{lemma}[theorem]{Lemma}}{}
\makeatother
\newcommand{\sn}{{\mathbb S}^{n-1}}
\newcommand{\tH}{\mathcal{H}}
\newcommand{\tk}{\mathcal{K}}
\title[Anisotropic fractional area measures]{Anisotropic fractional area measures}
\author{Xiaxing Cai}
\date{Source: arXiv:2510.05279v1 (CC BY 4.0)}
\begin{document}
\maketitle

\noindent\emph{Source and licence.} Anisotropic fractional area measures. Version arXiv:2510.05279v1 (CC BY 4.0), distributed under the Creative Commons Attribution 4.0 International licence, as recorded in the arXiv licence metadata for this exact version and shown on the abstract page https://arxiv.org/abs/2510.05279v1. Full rights notice: licenses/arxiv-2510.05279-CC-BY-4.0.txt.\par\smallskip

\noindent\textit{Editorial scope.} Counted unit: the Lemma whose source label is \texttt{conv} in the article (source0.tex lines 589--591, source label \texttt{conv}), delivered together with the article's own complete original proof of it (source0.tex lines 593--670, one \texttt{proof} environment of 78 lines). No mathematics is written, repaired or re-derived: statement and proof are the article's own text line for line. Required context retained uncounted: none. Every definition, display and label of those lines is retained unchanged. Omitted from this package and not redistributed: 1--588, 592--592, 671--1135. Nothing inside a retained excerpt is omitted or altered: every retained body is delivered exactly as the article's lines are written, so no omission receipt of any kind accompanies this package. Citations: the retained excerpts carry no citation command at all, so no bibliography entry of the article is transcribed and no reference is redistributed. Reference adaptation: none was needed and none was made; every reference inside the delivered unit resolves inside it, so no unresolved reference or citation is produced. Printed slips kept literal: no printed word, symbol, hypothesis or number of the counted statement or of its proof was altered. Layout and class: the article's own class is \texttt{amsart} with packages including amsmath, amsfonts, amssymb, eucal, mathrsfs, latexsym, enumitem, bbm, mathtools, exscale, cases, hyperref; the wrapper is the lane's pinned \texttt{amsart} editorial wrapper at 10pt on A4, which restates the theorem-like environments the retained text needs and adds the article's own macro definitions that the retained text needs (\texttt{\textbackslash sn}, \texttt{\textbackslash tH}, \texttt{\textbackslash tk}), with extra wrapper packages (bm, bbm, dsfont, xspace, cancel, mathtools). The delivered numbering is the unit's own. Assets: no retained span contains a figure environment, a graphics-inclusion command, a PDF-page inclusion, a table or a TikZ picture; no image file is copied into this package, the package declares no asset dependency and no image file is redistributed with it. Licence: version arXiv:2510.05279v1 of this article is distributed under the Creative Commons Attribution 4.0 International licence, as recorded in the arXiv licence metadata for this exact version and shown on the abstract page https://arxiv.org/abs/2510.05279v1. A full rights notice is delivered at licenses/arxiv-2510.05279-CC-BY-4.0.txt. Characters: every retained excerpt is pure ASCII, so the delivered engine, XeLaTeX with its default Latin Modern text font, reports no missing character.
\begin{lemma}\label{conv} Suppose $K\in\tk^n$ has $C^2$ boundary and everywhere positive Gauss curvature. Let $L\in\tk_e^n$ and $z\in\partial K$. Then
    \[\lim_{s\rightarrow 1^-}(1-s)\int_{\mathbb S_z^+}\rho_L(u)^{n+ s}X_K(z,u)^{-s}du=\int_{\sn\cap \nu_K(z)^{\perp}}\rho_L(\theta)^{n+1}\kappa_K(z,\theta)d\tH^{n-2}( \theta).\]
\end{lemma}

\begin{proof}
    Since $X_K(z,u)$ and $\kappa_K(z,\theta)$ are invariant under the translation of $z$, we assume $z=o$ in what follows.    
    Let $e_n=-\nu_K(o)$ and let $e_1,\ldots ,e_{n-1}$ be the eigenvectors of $\nabla^2h_K(-e_n)$, with  corresponding eigenvalues $1/\kappa_1,\ldots, 1/\kappa_{n-1}$. The basis $e_1,\ldots, e_n$ is chosen to be orthonormal, since $\nabla^2h_K(-e_n)$ is self-adjoint.
    
    In a small neighborhood of $o$, the boundary $\partial K$ can then be locally represented as the graph of a function:
    \begin{equation}\label{loc}
        y_n=f(y^{\prime})=\frac{1}{2}(\kappa_1y_1^2+\cdots +\kappa_{n-1}y_{n-1}^2)+o(|y^\prime|^2),
    \end{equation}
    where $y=(y^{\prime},y_n)=(y_1,\ldots,y_{n-1},y_n)$ is the coordinate with respect to the orthonormal frame $(o; e_1,\ldots,e_n)$.

    In this coordinate system, $X_K(o,u)=\rho_K(u)$ is positive if and only if $ u\cdot e_n>0$. We denote the corresponding integration domain by $D=\{u\in\sn:  u\cdot e_n>0\}$. 
    
    Let $\epsilon>0$ be arbitrary. Suppose $y^{\prime}=r\theta$, where $\theta\in \mathbb S^{n-2}$. Note that  $f(r\theta)=O(r^2)$, when $r>0$ is sufficiently small. It follows that
    %\begin{align*}
    %    \left|\frac{(r\theta, f(r\theta))}{\sqrt{r^2+f(r\theta)^2}}-(\theta,0)\right|^2=2-\frac{2r}{\sqrt{r^2+f(r\theta)^2}}=2-\frac{2}{\sqrt{1+f(r\theta)^2/r^2}}\rightarrow 0,
    %\end{align*}

    \[
    \left|\frac{(r\theta, f(r\theta))}{\sqrt{r^2+f(r\theta)^2}} - (\theta, 0)\right|^2
    = 2 - \frac{2}{\sqrt{1 + f(r\theta)^2 / r^2}} \to 0,
    \]
     as $r\rightarrow 0^+$. Therefore we choose $\delta>0$ such that \eqref{loc} holds for $|y^{\prime}|<\delta$ and 
    \begin{equation}\label{lem 1 appro 1}
        \left|\rho_L\bigg(\frac{(r\theta, f(r\theta))}{\sqrt{r^2+f(r\theta)^2}}\bigg)^{n+1}-\rho_L(\theta,0)^{n+1}\right|<\epsilon,
    \end{equation}
    when $r\in(0,\delta)$. We denote by
    \[D_1=\{u\in D: |y^{\prime}|\in (0,\delta), \text{where}~~(y^{\prime},f(y^{\prime}))=\rho_K(u)u\in\partial K \}.\]
    
    %Let $d\tH^{n-1}(y)$ be the surface area element of $\partial K$ and let $d\tH^{n-1}(y^{\prime})$ be the volume element in $e_n^{\perp}$. We have

    Recall that
    \begin{equation}\label{coor trans1}
        \nu_K(y)=\frac{(\nabla f(y^{\prime}),-1)}{\sqrt{1+|\nabla f(y^{\prime})|^2}},\quad d\tH^{n-1}_{\partial K}(y)=\sqrt{1+|\nabla f(y^{\prime})|^2}d\tH^{n-1}_{e_n^{\perp}}(y^{\prime}).
    \end{equation}
    For the spherical Lebesgue measure $du$ and $y^{\prime}\in e_n^{\perp}$ with $|y^{\prime}|\in(0,\delta)$, we also have
    \begin{equation*}
        \rho_K(u)^ndu=| y\cdot \nu_K(y) |d\tH^{n-1}(y),\quad y^{\prime}\cdot  \nabla f(y^{\prime})=2f(y^{\prime})+o(|y^{\prime}|^2).
    \end{equation*}
    Combining with \eqref{coor trans1}, we obtain
    \begin{equation}\label{coor trans2}
        \rho_K(u)^ndu=(f(y^{\prime})+o(|y^{\prime}|^2))d\tH^{n-1}(y^{\prime})
    \end{equation}

    We first consider the integral on $D_1$. By \eqref{coor trans2} and the coordinate transform $\rho_K(u)u=(y^{\prime}, f(y^{\prime}))$ , we have
    \begin{align*}
        (1-&s)\int_{D_1}\rho_L(u)^{n+s}\rho_K(u)^{-s}du\\
        &=(1-s)\int_{\{y^{\prime}\in e_n^{\perp}:~|y^{\prime}|<\delta\}}\rho_L\Big(\frac{(y^{\prime},f(y^{\prime}))}{|(y^{\prime},f(y^{\prime}))|}\Big)^{n+s}|(y^{\prime},f(y^{\prime}))|^{-s-n}(f(y^{\prime})+o(|y^{\prime}|^2))d\tH^{n-1}(y^{\prime})\\
        &=(1-s)\int_{\mathbb S^{n-2}}\int_0^{\delta}\rho_L\Big(\frac{(r\theta,f(r\theta))}{|(r\theta,f(r\theta))|}\Big)^{n+s}|(r\theta,f(r\theta))|^{-s-n}(f(r\theta)+o(r^2))r^{n-2}drd\tH^{n-2}(\theta).
    \end{align*}
    Note that $f(r\theta)=\frac{1}{2}r^2\kappa_K(o,\theta)+o(r^2)$. Then
    \begin{align*}
        &(1-s)\int_{D_1}\rho_L(u)^{n+s}\rho_K(u)^{-s}du\\
        \leq &(1-s)\int_{\mathbb S^{n-2}}\int_0^{\delta} \rho_L\Big(\frac{(r\theta,f(r\theta))}{|(r\theta,f(r\theta))|}\Big)^{n+s}\Big(1+\frac{r^2\kappa_K(o,\theta)^2}{4}+o(r^2)\Big)^{\frac{-s-n}{2}}r^{-s}\\ 
        & \quad\quad\quad\quad\quad\quad\quad\quad\quad\quad\quad\quad\quad\quad\quad\quad\quad\quad\quad\quad\quad\quad\quad\quad\quad\quad\cdot\Big(\frac{\kappa_K(o,\theta)}{2}+o(1)\Big)drd\tH^{n-2}(\theta)\\
        \leq &(1-s)\int_{\mathbb S^{n-2}}\int_0^{\delta}\rho_L\Big(\frac{(r\theta,f(r\theta))}{|(r\theta,f(r\theta))|}\Big)^{n+s}\frac{\kappa_K(o,\theta)}{2}r^{-s}drd\tH^{n-2}(\theta)+(1-s)o(1)\int
        _0^{\delta}r^{-s}dr.
    \end{align*}
    By \eqref{lem 1 appro 1}, we have
    \begin{align*}
        (1-s)\int_{D_1}&\rho_L(u)^{n+s}\rho_K(u)^{-s}du\\
        \leq &(1-s)\int_{\mathbb S^{n-2}}\int_0^{\delta}(\rho_L(\theta,0)^{n+1}+\epsilon)\frac{\kappa_K(o,\theta)}{2}r^{-s}drd\tH^{n-2}(\theta)+o(1)\delta^{1-s}\\
        =&\frac{\delta^{1-s}}{2}\int_{\mathbb S^{n-2}}(\rho_L(\theta,0)^{n+1}+\epsilon)\kappa_K(o,\theta)d\tH^{n-2}(\theta)+o(1)\delta^{1-s}.
    \end{align*}
    Let $s\rightarrow 1^{-}$. By the arbitrariness of $\epsilon$, we have
    \[\limsup_{s\rightarrow1^-}(1-s)\int_{D_1}\rho_L(u)^{n+s}\rho_K(u)^{-s}du\leq \frac{1}{2}\int_{\mathbb S^{n-2}}\rho_L(\theta,0)^{n+1}\kappa_K(o,\theta)d\tH^{n-2}(\theta)+o(1).\]

    For $u\in D\backslash D_1$ and $(y^{\prime},f(y^{\prime}))=\rho_K(u)u$, we have $\rho_K(u)\ge \delta $. Therefore $\rho_L(u)^{n+s}\rho_K(u)^{-s}$ is uniformly bounded, i.e., $\rho_L(u)^{n+s}\rho_K(u)^{-s}\leq c$ for some constant $c$ and for all $u\in D\backslash D_1$ and $s\in(0,1)$. Then
    \[\limsup_{s\rightarrow 1^{-}}(1-s)\int_{D\backslash D_1}\rho_L(u)^{n+s}\rho_K(u)^{-s}du\leq \limsup_{s\rightarrow 1^{-}}(1-s)\int_{D\backslash D_1} c du=0,\]
    which implies
    \[\limsup_{s\rightarrow 1^-}(1-s)\int_D\rho_L(u)^{n+s}\rho_K(u)^{-s}du\leq \frac{1}{2}\int_{\mathbb S^{n-2}}\rho_L(\theta,0)^{n+1}\kappa_K(o,\theta)d\tH^{n-2}(\theta)
    +o(1).\]
    Letting $\delta\rightarrow 0^+$, we then have
    \[\limsup_{s\rightarrow 1^-}(1-s)\int_D\rho_L(u)^{n+s}\rho_K(u)^{-s}du\leq \frac{1}{2}\int_{S^{n-2}}\rho_L(\theta,0)^{n+1}\kappa_K(o,\theta)d\tH^{n-2}(\theta).\]

    By the other side of the bound in \eqref{lem 1 appro 1} and Fatou's Lemma, we repeat the same arguments and get
    \[\liminf_{s\rightarrow 1^-}(1-s)\int_D\rho_L(u)^{n+s}\rho_K(u)^{-s}du\geq \frac{1}{2}\int_{S^{n-2}}\rho_L(\theta,0)^{n+1}\kappa_K(o,\theta)d\tH^{n-2}(\theta),\]
    which concludes the proof.
\end{proof}

\end{document}
